\chapter{Zakończenie}

Celem pracy było sprawdzenie, czy mały model językowy dostrojony za pomocą LoRA
może klasyfikować wiadomości jako \texttt{ham} albo \texttt{spam} oraz jak
wypada w porównaniu z prostszymi klasyfikatorami tekstu. Przygotowano wspólny
zbiór danych, porównano cztery sposoby sformułowania zadania dla Qwen3-0.6B i
oceniono wybrane modele na wiadomościach pochodzących z czterech odrębnych
źródeł.

W badanych warunkach dostrojony Qwen3-0.6B osiągnął wyniki wyższe od metod
bazowych. Najlepsze rezultaty uzyskały generowanie odpowiedzi
ustrukturyzowanej i ocena następnego tokenu, przy czym zaobserwowana różnica
mieściła się w granicach niepewności pomiaru. Do dalszych porównań wybrano
ocenę następnego tokenu, która nie wymaga generowania pełnej odpowiedzi ani jej
parsowania. Najwyższe wyniki walidacyjne często pojawiały się przed końcem
treningu.

Wyższa skuteczność Qwen3 wiązała się z większym kosztem inferencji niż w
przypadku DistilBERT i klasyfikatorów opartych na TF-IDF. W środowisku
produkcyjnym model językowy mógłby uzupełniać reguły, reputację nadawcy i
tańszy klasyfikator tekstu, szczególnie dla wiadomości trudnych do
jednoznacznego rozstrzygnięcia. Ocena takiego rozwiązania wymaga jednak danych
z bieżącego ruchu pocztowego i pomiarów pod obciążeniem. W pracy przygotowano i
oceniono klasyfikator tekstowy oparty na małym modelu z otwartymi wagami.
